%% =====================================================================
%%  Écouter les audios
%% =====================================================================
\liminaire{audios}{Écouter les audios}{Plus de 250 pistes, gratuites, sur votre téléphone ou votre ordinateur}

On n'apprend pas une langue avec les yeux : chaque dialogue, chaque liste de mots et chaque
exercice d'écoute de ce livre est enregistré. Il n'y a rien à installer.\par\vspace{5mm}

\noindent\begin{tikzpicture}
  \foreach \n/\ico/\titre/\txt in {1/micro/Ouvrez l'appareil photo/de votre téléphone (ou une application de lecture de QR codes).,
      2/cible/Visez le QR code/imprimé à côté de l'activité\, comme si vous preniez une photo.,
      3/casque/Touchez le lien qui apparaît/: la piste se lance. Mettez en pause\, revenez en arrière\, réécoutez.}{%
    \pgfmathsetmacro{\x}{(\n-1)*58}
    \fill[otNavyTint,rounded corners=2mm] (\x mm,0) rectangle ++(54mm,40mm);
    \fill[otNavy] (\x mm+9mm,31mm) circle (5mm);
    \node[text=white,font=\otBold\fontsize{11}{12}\selectfont] at (\x mm+9mm,31mm) {\n};
    \otDessinIcone{\ico}{otNavy!35}{otNavyTint}{9mm}{(\x mm+45mm,31mm)}
    \node[anchor=north west,text width=46mm,inner sep=0pt,align=left] at (\x mm+4mm,23mm)
      {{\otSemi\fontsize{10}{12}\selectfont \titre}\par\vspace{1mm}{\fontsize{9.6}{12}\selectfont\txt}};}
\end{tikzpicture}\par\vspace{6mm}

\begin{paire}
\begin{tcolorbox}[otBase,colback=white,colframe=otRule,boxrule=0.5pt]
\otTitreBoite{casque}{otNavy}{Le numéro de piste}{}
\begin{center}\otPastillePiste{1.05}\end{center}
{\fontsize{10}{12.6}\selectfont Piste \textbf{1.05} : unité \textbf{1}, piste \textbf{5}, suivie de sa durée.
Les dialogues ont deux pistes : une \textbf{version lente} pour bien distinguer chaque mot, et une
\textbf{version normale}, au rythme d'une vraie conversation.\par}
\end{tcolorbox}
\begin{tcolorbox}[otBase,colback=otNavyTint,colframe=otNavyTint,halign=center]
\otFixeURL{index.html}%
{\otSemi\fontsize{7}{8.5}\selectfont\color{otNavy}\otCaps[11]{Tous les audios du livre}}\par\vspace{2mm}
{\color{otInk}\otQR{25mm}{}}\par\vspace{2mm}
{\otMedium\fontsize{9.4}{11.4}\selectfont\otAudioAffiche}\par\vspace{1mm}
{\fontsize{9.2}{11.4}\selectfont\color{otMuted}Sur ordinateur : toutes les pistes, unité par unité, en lecture ou en téléchargement (MP3).\par}
\end{tcolorbox}
\end{paire}

\otIntertitre{Les voix que vous allez entendre}
Les pistes sont enregistrées en \textbf{anglais américain standard}, l'accent le plus entendu dans
le monde (films, séries, Internet). Emma, la collègue londonienne de Léa, parle avec un
\textbf{accent britannique} : vous apprendrez ainsi à comprendre les deux. Les enregistrements sont
réalisés avec des voix de synthèse neuronales de haute qualité, toutes contrôlées une à une.

\otIntertitre{Trois façons d'écouter}
\begin{tcbraster}[raster columns=3,raster equal height=rows,raster column skip=3.5mm,raster before skip=1mm]
\otCartePrg{casque}{otNavy}{Écoute globale}{Écoutez sans lire et cherchez à comprendre la situation : qui parle, où, pourquoi ?}
\otCartePrg{livre}{otGreenDark}{Écoute active}{Réécoutez en suivant le texte, phrase par phrase. Mettez en pause dès qu'un mot vous échappe.}
\otCartePrg{micro}{otOrangeDark}{Écoute « en écho »}{Répétez juste après la voix, en imitant son rythme et sa mélodie : c'est la technique du \en{shadowing}.}
\end{tcbraster}
